%!TEX program = lualatex
\documentclass[ngerman,headheight=70pt]{scrartcl}
\usepackage[ngerman]{babel}
\usepackage{blindtext}
\usepackage[utf8]{luainputenc}
\usepackage{scrlayer-scrpage}
\usepackage[T1]{fontenc}
\usepackage{fontspec}
\usepackage{multicol}
\usepackage{enumitem}
\usepackage{epigraph}
\usepackage[german=quotes]{csquotes}
\usepackage[bottom=2cm,footskip=8mm]{geometry}
\setmainfont{Source Sans Pro}
\MakeOuterQuote{"}

\setlength{\parindent}{0pt}
\setlength{\parskip}{6pt}

\pagestyle{plain.scrheadings}
\thispagestyle{scrheadings}
\lohead{\vspace*{-1.5cm}\includegraphics[scale=0.4]{../images/stupa_logo_rot.pdf}}
\rohead{\\Präsidium des\\Studierendenparlaments}

\lofoot[Universität Hamburg · Präsidium des Studierendenparlaments\\\hrule Von-Melle-Park 5 · D-20146 Hamburg]{Universität Hamburg · Präsidium des Studierendenparlaments\\\hrule Von-Melle-Park 5 · D-20146 Hamburg}
\cofoot[]{}
\cohead[\thepage]{}

\renewcommand{\thesubsection}{TOP \arabic{subsection}}
\renewcommand{\thesubsubsection}{\arabic{subsubsection}.}
\setcounter{subsection}{-1}

\begin{document}
    UHH · StuPa-Präsidium · Von-Melle-Park 5 · D-20146 Hamburg

    \section*{Beschluss des Studierendenparlamentes vom 10. November 2016}
    \subsection*{Die Kampagne der Verfassten Studierendenschaft
    Kein G20, nirgends: Für eine bessere Welt!}

    Die Stadt Hamburg und die Bundesregierung wollen den kommenden Gipfel der
    20 größten Industrie- und Schwellenländer am 7. und 8. Juli 2017 in Hamburg
    stattfinden lassen. Die Verfasste Studierendenschaft beteiligt sich
    daraufhin in verschiedenen Formen -- durch Aufklärung und öffentlichen Protest,
    wie Erstellung von Informationsmaterialien, Planung eines Alternativ-Gipfels,
    Demonstrationen -- an den stadtweiten Protesten gegen den Wahnsinn von G20
    und für etwas Besseres.

    \blockquote[Stadt Hamburg: Das G20-Gipfeltreffen in Hamburg, http://www.hamburg.de/g20-gipfel/hintergrund]{
    Hamburg -will im Geiste des Friedens eine Mittlerin zwischen
    allen Erdteilen und Völkern der Welt sein-, heißt es in der Präambel unserer
    Landesverfassung. Dieses Anliegen lässt sich beim G20-Gipfel verwirklichen.}

    Dass der Hamburgische Senat zum G20-Gipfeltreffen unter Legitimationsdruck
    steht, zeigt das Zitieren der Präambel der Hamburger Verfassung in ihrer
    offiziellen Meldung: Hier wird suggeriert, bei G20 ginge es um so etwas wie
    Völkerverständigung. Doch was bei Olympia in Hamburg schon ohne Brot
    stattfinden sollte, wird jetzt auch ohne Spiele versucht durchzusetzen:
    Militarisierung, Einschränkungen der Grundrechte, Überwachung in Hamburg und
    Armut, Krieg, Austeritätspolitik und Umweltzerstörung für alle.

    Wie schon in der Kampagne um das NOlympia-Bündnis schließt sich die Verfasste
    Studierendenschaft im Rahmen des universitären Leitbildes, in dem die die
    Mitglieder der Universität dazu aufgefordert werden, "zur Entwicklung einer
    humanen, demokratischen und gerechten Gesellschaft" beizutragen, den Protesten
    an und setzt sich das Ziel, einen G20-Gipfel in der Stadt zu verhindern. Kein
    G20 in Hamburg oder irgendwo sonst! Stattdessen: Wir treten für ein gutes
    Leben für alle ein, das heißt konkret für Hamburg: Ende der
    Schuldenbremsenpolitik, Ausfinanzierung der Universitäten zum Allgemeinwohl,
    Stopp der Waffenexporte über den Hamburger Hafen und ein Bleiberecht für alle!

    \subsubsection*{Die Gruppe der 20 ökonomisch mächtigsten  Industrie- und
    Schwellenländer}

    Nach der Auflösung der Sowjetunion 1991, mit der auch der Warschauer Pakt zu
    Ende ging, wurde in den 1990er Jahren der globale Kapitalismus verfestigt,
    mit dem Wohlstand und Frieden für alle propagiert wurde. Stichwort ist hier
    das von dem Politologen Francis Fukuyama propagierte "Ende der Geschichte":
    Große Richtungsentscheidungen würden der Vergangenheit angehören, die ganze
    Welt entwickele sich nun sozial-marktwirtschaftlich usw. nach westlichem
    Vorbild. Ab Mitte der 90er zeigte sich allerdings bereits deutlich die
    Krisenhaftigkeit einen solches Kapitalismus in diversen Finanzkrisen. Mit
    der sogenannten Tequila-Krise (Mexiko) Mitte der 90er, der Russland-Krise
    sowie der Asienkrise Ende der 90er wurden schlussendlich alle Illusionen
    zerstört. Aus den Krisen heraus resultierte eine neoliberale Neuordnung, die
    mit einer verstärkten Liberalisierung die kapitalistisch-neoliberale
    Krisenzeit beenden sollte, statt mit umfassenden Sozialreformen. Die
    Liberalisierung fand Ausdruck in dem globalen Kampf um Ressourcen, erkennbar
    in der neokolonialistischen Entwicklungspolitik der ökonomisch betrachtet
    großen Staaten, aber auch im Abbau der Sozialsysteme dort auf nationaler
    Ebene. Die G20 sind dabei ein zentraler Akteur und finden ihren Ursprung in
    der Zeit nach der Asienkrise 1997/98. 1999 trafen sich die Vertreter der 19
    mächtigsten Industrienationen und Schwellenländer im Rahmen eines
    Finanzministertreffens in Berlin, um globale finanzpolitische Maßnahmen zur
    Stabilisierung der Weltwirtschaft zu diskutieren.

    Die G20 sind dabei stets ein rein informelles Bündnis, dessen Beschlüsse somit
    keine rechtliche Bindung haben. Die G20 legitimieren sich nur dadurch, dass
    sie großteils aus demokratisch gewählten Regierungen bestehen. Neben der
    Selbstlegitimierung fehlt also die Legitimierung durch die Weltgemeinschaft,
    in der die G20 agieren.

    \subsubsection*{Die G20: Nicht der Bekämpfer der globalen Krisen, sondern ein
    zentraler Auslöser dieser}

    Wer sich vor allem mit Appellen an die G20 wendet und sich so eine globale
    Politik der Nachhaltigkeit, sozialer Gerechtigkeit und friedlicher
    Konfliktlösung erhofft, verkennt die Realität: Der Kapitalismus und seine
    Hauptakteure, zu denen die G20 als Regierungen besonders wirtschaftsstarker
    Nationen gehören, verschärfen die globale Situation immer weiter, was sich
    politisch heute u.a. in der Polarisierung der Gesellschaft widerspiegelt.

    Die grundsätzliche Widersprüchlichkeit zwischen Kapital und Arbeit, die sich
    in den kapitalistischen Krisen immer offener zeigt, wird durch eine globale
    Privatisierungs- und Austeritätspolitik vermeintlich beantwortet. Hierbei
    spielen globale Institutionen wie der Internationale Währungsfond, die World
    Trade Organisation und die Weltbank (z.B. als Kreditgeber), eine wesentliche
    Rolle. Diese Institutionen sind dabei eben nicht unabhängig, sondern
    unterliegen direkt dem Einfluss der mächtigsten Wirtschaftsnationen. Ziel
    dieser Politik ist nicht die Stabilisierung von "Entwicklungsländern" oder
    "Pleitestaaten" wie Griechenland, sondern vor allem eine Stärkung der
    nationalen Privatwirtschaft der agierenden Regierungen. Die Folgen sind der
    Abbau von Sozialsystemen und die Ausbeutung der Bevölkerung in den betroffenen
    Staaten, die in Menschenrechtseinschneidungen enden und so zu Verfolgung und
    Flucht führen.

    Darüber hinaus befinden sich unter den G20 die größten Waffen- und
    Rüstungshersteller. Hier zeigt sich, wie paradox die Hoffnung ist, die G20
    würden zu globalen Frieden beitragen. Stattdessen treten sie als Aggressor
    auf und beteiligen sich direkt oder indirekt in vielen Kriegen wie dem Irak-,
    dem Afghanistan- und dem Libyenkrieg. Dabei fungieren Kriege auch als
    Strategie eines imperialistischen und neokolonialistischen Politikstils, bei
    dem die Sicherung großer Rohstoffvorkommen und die Monopolbildung im
    Mittelpunkt steht.

    Auch auf die ökologische Frage kann nicht mit einer Antwort der G20 gerechnet
    werden. Die massive Durchsetzung imperialistischer Kapitalinteressen
    innerhalb der global herrschenden Wachstumsideologie führt zur
    unkontrollierten Zerstörung der Natur. Die umfassende Rodung von Regenwald,
    die industrielle Überfischung und der wachsende CO2-Ausstoß sind nur einige
    Beispiele dafür, dass im Rahmen der aktuellen neoliberalen Politik keine
    bessere Aussicht für die Biodiversität und das Klima zu erwarten ist.

    \subsubsection*{Der G20-Gipfel: Reproduktion der Kapitalinteressen auf lokaler
    Ebene}

    Ein künstlicher See in einer Veranstaltungshalle, der lediglich für ein Foto
    der beteiligten Regierungschefs dient -- das ist nur ein Ausdruck des
    Wahnsinns, den ein Austragungsort eines G20-Gipfels erwarten kann. Der
    jährliche Gipfel dient vor allem  als  Medienevent, um den Menschen
    vorzugaukeln, es würde ein tatsächliches Krisenmanagement stattfinden, denn:
    In Wirklichkeit sind die Staats- und Regierungschefs der G20 nicht in der
    Lage die globalen Krisen zu lösen. Ihre Politik ist gescheitert, da sie nicht
    anerkennen, dass die Krisenhaftigkeit dem Kapitalismus selbst innewohnt und
    sie die Krisen wiederum benötigen, um dieses System am Leben zu erhalten.

    G20-Gipfel, wie der 2010 in Toronto, beweisen außerdem, wie auch
    innenpolitisch durch die Regierung des Gastgeberlandes eine Militarisierung
    und verstärkte Überwachung durchgesetzt wird: Um die Sicherheitskriterien
    der Geheimdienste der Mitgliedsstaaten zu erfüllen, wird ein Sicherheitskonzept
    mit Kosten im zehnstelligen Bereich aufgestellt, das die Grundrechte der
    Zivilbevölkerung enorm einschränkt. Hierbei kam es in Toronto zur Einführung
    einer flächendeckenden Kameraüberwachung im Bereich der Innenstadt und mit
    der Einführung eines Sonderpolizeirechts zu über 900 festgenommen
    Demonstrant*innen, was den größten Massenarrest in der kanadischen Geschichte
    darstellte.

    Für Hamburg bedeutet die Durchführung des G20-Gipfels in den Messehallen und
    somit im Zentrum zwischen Universität, Schanzenviertel und St. Pauli, also
    ein massives Polizeiaufkommen und eine mögliche erneute Diskussion über einen
    Bundeswehreinsatz im Inneren, eine Ausweitung und Wiedereinführung von
    Gefahrengebieten nur unter anderem Namen, mit hoher Wahrscheinlichkeit ein
    umfangreiches öffentliches Kameraüberwachungssystem und nicht zuletzt hierfür
    Kosten möglicherweise im Milliardenbereich, die in einem enormen Widerspruch
    zur aktuellen "Sparpolitik" des Senats sowie der Bundesregierung stehen.

    \subsubsection*{Die Verfasste Studierendenschafft muss sich den G20 gegenüberstellen
    und für eine bessere Welt einstehen, weil sie so auch Wissenschaft versteht.}

    \blockquote[Naomi Klein, Foreword: G20 Trials and the War on Activism. In:
    T. Malleson/D. Wachsmuth (Hgg.): Whose Streets? The Toronto G20 and the
    Challenges of Summit Protest. Toronto 2011. S. xi-xvii.]{
    These are fights we can if we build coalitions like the ones we saw on the
    streets of Toronto during the G20. (...) Our Government fears those
    coalitions, fears the prospect of a truly mass social movement, and we can
    see that fear in arrest and prosecution patterns. (...) And when we refuse
    to forget what happend here during the G20, when we demand accountability
    for the real criminals, and the	freedom of our friends, we are fighting not
    just for the past but for the future.}

    Die ermutigenden Worte von Naomi Klein nach den umfangreichen Repressionen
    gegen die G20-Gipfelgegner*innen von Toronto 2010 zeigen, dass die Regierungen
    Angst vor sozialen, kapitalismuskritischen Protesten haben und in diesen also
    das Potenzial für eine gesellschaftliche Transformation liegt.

    Aus den Protesterfahrungen in Toronto zeigt sich allerdings auch, dass der
    Gipfel in Hamburg und überall sonst zu verhindern ist. Denn: Die Repression
    wird auch hier all die Gipfelgegner*Innen treffen und die Stadt wird auch
    hier einer neuen Form der Überwachung ausgesetzt und diese reale Gefahr kann
    nur durch das vorzeitige Aus des "Image-Treffens" verhindert werden. Ein
    Ansatz der dagegen von Beginn an nur auf die umfangreiche Protestkultur
    während des Gipfels setzt, verkennt diese Gefahr. Eine Verhinderung des Gipfels
    wäre eine grundsätzliche Absage, eben auch an die Inhalte, die von den G20
    umgesetzt werden und wäre der größte kapitalismuskritische Erfolg der
    Protestbewegung.

    Dafür ist es erforderlich bei allen Aktivitäten -- einschließlich den hier
    zum Schluss des Antrags genannten -- zu erarbeiten und deutlich zu machen,
    was die fortschrittliche, friedenspolitische und soziale Alternative zu
    den von den G20 repräsentierten Positionen ist. Diese Alternative muss
    durch die Bevölkerungen international erwirkt werden.

    Die VS beteiligt sich deshalb mannigfaltig:
    \begin{itemize}
        \item Planung einer Info/Diskussionsveranstaltung für den November, um
              die wissenschaftliche Auseinandersetzung an der Uni zu beginnen.
        \item Erstellung einer Aufklärungsbroschüre, um die Bevölkerung der Stadt
              Hamburg zu informieren und zu Widerstand zu ermutigen.
	    \item Planung einer Demonstration für Anfang 2017, um die die mehrheitliche
              Ablehnung der Bevölkerung auf die Straße zu bringen.
        \item Beteiligung am Alternativgipfel, der aufzeigt, wieso die G20 keine
              Lösungen für die globalen Krisen bereithalten, und Alternativen zum
              jetzigen Neoliberalismus bietet. Nutzung der Uni-Infrastruktur
              für die Durchführung des Alternativgipfels.
    \end{itemize}
\end{document}
